\documentclass[12pt]{letter}
\usepackage[margin=1in,letterpaper]{geometry}

\signature{Alexei Drummond, University of Auckland\\[0.5cm]Alex Popinga, Colorado State University}

\begin{document}

\begin{letter}{Editor\\Journal of Quantitative Analysis in Sports}

\opening{Dear Professor Glickman,}

We are pleased to submit our manuscript, ``Bayesian inference of the climbing grade scale,'' for consideration in the Journal of Quantitative Analysis in Sports.

In this paper, we apply the dynamic Bradley-Terry model to rock climbing ascent data from a large public logbook (thecrag.com) to quantify the fundamental scale parameter of climbing grade systems. Using Bayesian MCMC inference implemented in Stan, we estimate that each increment in the Ewbank, French and UIAA sport climbing grade systems corresponds to approximately a doubling in difficulty, while the V-grade bouldering scale corresponds to approximately a tripling. We show that climbing grades behave as a logarithmic scale of difficulty, analogous to decibels or stellar magnitude, and discuss connections to the Weber-Fechner law of psychophysics.

Rock climbing has been an Olympic sport since the 2021 Tokyo Games, and bouldering and climbing gyms have exploded in popularity over the last decade. We believe this work is well suited to JQAS, as it brings rigorous statistical methodology to bear on a question of direct interest to this rapidly growing sport and the wider sports analytics community. A preprint of this work is available on arXiv (2111.08140).

We suggest Nick Draper (University of Canterbury) and Ioannis Ntzoufras (Athens University of Economics and Business) as potential reviewers, given their respective expertise in quantitative research on climbing grades and Bayesian models in sports.

All authors have approved this submission.

Thank you for your consideration.

\closing{Sincerely,}

\end{letter}
\end{document}
